Industrial hoisting systems require not only post-event fault recognition, but also short-horizon prediction of degradation states before a fault process becomes critical. In the hoister overspeed dataset studied here, the key challenge is not simply average classification accuracy: standard time-series classifiers can achieve high accuracy while completely missing the rare second-level degradation state. We formulate the task as fixed-horizon future-state classification, where a current multivariate sensor window is used to predict the state label one step ahead. To address rare-state collapse around transition boundaries, we propose SGTONetV6, a dual-mode model that decouples conservative future-state classification from boundary-constrained rare-fault triggering. The model uses a patch temporal encoder, a conservative multiclass classifier, a patch-attentive rare context module, and a calibrated rare-trigger override constrained by boundary and precursor-state semantics. On a private hoister dataset with 27 files and 37,417 timestamps, SGTONetV6 achieves the best macro-F1 among the tested short-horizon models and recovers the rare class-9 state with F1 0.5556, while DLinear, TimesNet, iTransformer, PatchTST, and a conservative SGTO baseline all obtain class-9 F1 of 0.0000. Ablations show that the rare trigger, boundary constraint, fallback prior, and patch-attentive rare context are all necessary for reliable rare-state recovery.
